\documentclass[10pt,a4paper]{amsart}
\usepackage[margin=19mm]{geometry}
\usepackage{amsmath,amssymb,mathtools}
\usepackage[scr=rsfs,cal=euler]{mathalfa}
\usepackage{xfrac}
\usepackage[unicode,hidelinks,bookmarks=false]{hyperref}
\newcommand{\ie}{i.e.\ }
\newcommand{\cf}{cf.\ }
\newcommand{\resp}{respectively\ }
\newcommand{\Subsection}{\subsection}
\providecommand{\nobreakdash}{\nobreak\hbox{-}}
\setlength{\emergencystretch}{2em}
\multlinegap0pt
\usepackage{bm}
\usepackage{bbm}
\usepackage{dsfont}
\usepackage{xspace}
\usepackage{cancel}
\usepackage{mathtools}
\usepackage{amsthm}
\makeatletter
\@ifundefined{theorem}{\newtheorem{theorem}{Theorem}}{}
\@ifundefined{proposition}{\newtheorem{proposition}[theorem]{Proposition}}{}
\makeatother
\title[An extension of Liebmann's Theorem to hypersurfaces with bou]{An extension of Liebmann's Theorem to hypersurfaces with boundary}
\author{Fl\'avio Fran\c{c}a Cruz, Barbara Nelli}
\date{Source: arXiv:2412.03368v3 (CC BY 4.0)}
\begin{document}
\maketitle

\noindent\emph{Source and licence.} An extension of Liebmann's Theorem to hypersurfaces with boundary. Version arXiv:2412.03368v3 (CC BY 4.0), distributed under the Creative Commons Attribution 4.0 International licence, as recorded in the arXiv licence metadata for this exact version and shown on the abstract page https://arxiv.org/abs/2412.03368v3. Full rights notice: licenses/arxiv-2412.03368-CC-BY-4.0.txt.\par\smallskip

\noindent\textit{Editorial scope.} Counted unit: the Proposition whose source label is \texttt{lema2} in the article (source0.tex lines 239--242, source label \texttt{lema2}), delivered together with the article's own complete original proof of it (source0.tex lines 245--295, one \texttt{proof} environment of 51 lines). No mathematics is written, repaired or re-derived: statement and proof are the article's own text line for line. Required context retained uncounted: teorema (lines 90--93); lema (lines 116--119). Every definition, display and label of those lines is retained unchanged. Omitted from this package and not redistributed: 1--89, 94--115, 120--238, 243--244, 296--450. Nothing inside a retained excerpt is omitted or altered: every retained body is delivered exactly as the article's lines are written, so no omission receipt of any kind accompanies this package. Citations: the retained excerpts carry no citation command at all, so no bibliography entry of the article is transcribed and no reference is redistributed. Reference adaptation: none was needed and none was made; every reference inside the delivered unit resolves inside it, so no unresolved reference or citation is produced. Printed slips kept literal: no printed word, symbol, hypothesis or number of the counted statement or of its proof was altered. Layout and class: the article's own class is \texttt{amsart} with packages including color, eucal, amssymb, url, amscd, amsfonts, latexsym, hyperref, xy, tikz; the wrapper is the lane's pinned \texttt{amsart} editorial wrapper at 10pt on A4, which restates the theorem-like environments the retained text needs and adds the article's own macro definitions that the retained text needs (none), with extra wrapper packages (bm, bbm, dsfont, xspace, cancel, mathtools). The delivered numbering is the unit's own. Assets: no retained span contains a figure environment, a graphics-inclusion command, a PDF-page inclusion, a table or a TikZ picture; no image file is copied into this package, the package declares no asset dependency and no image file is redistributed with it. Licence: version arXiv:2412.03368v3 of this article is distributed under the Creative Commons Attribution 4.0 International licence, as recorded in the arXiv licence metadata for this exact version and shown on the abstract page https://arxiv.org/abs/2412.03368v3. A full rights notice is delivered at licenses/arxiv-2412.03368-CC-BY-4.0.txt. Characters: every retained excerpt is pure ASCII, so the delivered engine, XeLaTeX with its default Latin Modern text font, reports no missing character.
\begin{theorem}
	\label{teorema}
	Let $\Sigma$ be a closed strictly convex  $(n-1)$-dimensional submanifold in a hyperplane $\Pi\subset \mathbb{R}^{n+1}$. Let $M$ be a connected compact embedded CMC hypersurface in $\mathbb{R}^{n+1}$ with boundary $\Sigma$. If $M$ is locally convex, then $M$ lies in one of the two halfspace defined by $\Pi$. In particular, if $\Sigma$ is a sphere, then $M$ is a $n$-ball or a spherical cap.
\end{theorem}

\begin{proposition}
	\label{lema}
	Let $M$ be as in Theorem \ref{teorema}.  Then, $M$ is locally strictly convex on its interior and its asymptotic directions on the boundary are necessarily tangent to $\Sigma$. 
\end{proposition}

\begin{proposition}
	\label{lema2}
	Let $M$ be as in Theorem \ref{teorema}.  Then, $M$, along $\Sigma$, is locally contained in one of the two halfspace of $\mathbb{R}^{n+1}$ determined by $\Pi.$ 
	\end{proposition}

\begin{proof}
 First we need to introduce some notations. The orientation of $M$ induces a natural orientation on its boundary as follows: we say that a basis $\{e_1, \ldots, e_{n-1}\}$ for $T_p\Sigma$ is positively oriented if $\{e_1, \ldots, e_{n-1}, \nu\}$ is a positively oriented basis for $T_pM$,  where $\nu$ denotes the inward pointing unit cornormal vector field along $\Sigma$. Let $\eta$  be the unitary vector field normal to $\Sigma$ in $\Pi$ which points outward with respect to the compact domain in $\Pi$ bounded by $\Sigma$. We denote by $\xi$ the unitary vector  normal to $\Pi$  in $\mathbb{R}^{n+1}$ which is compatible with $\eta$ and with the orientation of $\Sigma$, i. e., such that $\{\eta, e_1, \ldots, e_{n-1},\xi\}$ is a positively oriented basis of $\mathbb{R}^{n+1}$. We can assume that $\Pi$ is the horizontal plane that passes through the origin of $\mathbb{R}^{n+1}$.

Let $\{e_1, \ldots, e_{n-1}\}$ be (locally defined) a positively oriented frame field along $\Sigma$. Using this frame, we can write $\nu = N\times e_1\times\cdots\times e_{n-1} $ and $\eta= e_1\times\cdots\times e_{n-1} \times\xi$, since $\det( e_1, \dots, e_{n-1},\nu,  N)=\det( e_1, \dots, e_{n-1}, \eta, \xi)=1.$ From these expressions we compute
\begin{align*}
\eta =& e_1\times\cdots\times e_{n-1} \times \xi \\ =& e_1\times\cdots\times e_{n-1} \times (\langle\xi, N\rangle N+\langle \xi, \nu\rangle \nu)\\ =& -\langle\xi, N\rangle \nu+\langle \xi, \nu\rangle N.
\end{align*}
where   $\bar\nabla$ denotes  the usual connection of $\mathbb{R}^{n+1}$. 
Thus,
\begin{equation}
		\label{conta-angulos}
\langle \eta ,\nu\rangle=-\langle \xi, N\rangle \quad \textrm{and}\quad \langle \eta , N\rangle=\langle \xi, \nu\rangle.
\end{equation}
A direct computation gives that the second fundamental form of $M$ along of $\Sigma$ satisfies
\begin{equation}
	\label{conta-h-bordo}
h(e_i,e_j)=  h_{\Sigma} (e_i,e_j)\langle \xi, \nu\rangle,
\end{equation}
for any $1\leq i,j<n$, where $h_\Sigma$ denotes the second fundamental form of $\Sigma$ as a submanifold of $\Pi$. In fact,  as
\[
\bar \nabla_{e_i}e_j=\sum_k^{n-1}\langle \bar \nabla_{e_i}e_j, e_k\rangle e_k+\langle\bar \nabla_{e_i}e_j, \eta\rangle \eta+\langle\bar \nabla_{e_i}e_j, \xi\rangle \xi
\]
for any $1\leq i,j<n$, we have
\[
h(e_i, e_j)=\langle\bar \nabla_{e_i}e_j, \eta\rangle\langle\eta,N\rangle + \langle \bar \nabla_{e_i}e_j, \xi \rangle \langle\xi, N\rangle.
\]
It then follows from $\langle \bar\nabla_{e_i}e_j, \xi \rangle =0$ that
\[
h(e_i,e_j)=\langle \bar \nabla_{e_i}e_j, \eta\rangle  \langle N,\eta\rangle,
\]
then \eqref{conta-h-bordo} follows from \eqref{conta-angulos}.

Now we parameterize a neighborhood of $\Sigma$ in $M$ by setting  
\[
X(p,t)=\exp_{p}\left(t\nu(p)\right),
\]
$(p,t)\in \Sigma\times [0,\epsilon),$ for $\epsilon>0$ sufficiently small. Let $p\in \Sigma$ an arbitrary point of $\Sigma$. We claim that
 \[
 f(t):=\langle X(p,t), \xi \rangle >0
 \]
  for  $0<t<\epsilon$, if $\epsilon>0$ is small enough. Indeed, it follows from \eqref{conta-h-bordo} that 
  \[
  f'(0)=\langle  \nu, \xi\rangle= \frac{h(e_i, e_i)}{h_\Sigma(e_i, e_i)}(p).
  \] 
  for any $1\leq i<n$. Then, if $h(e_i,e_i)(p)>0$ for some $1\leq i<n$, we have that $f(0)=0$ and $f'(0)>0$, hence the claim is proved. On the other hand, if $h(e_i,e_i)(p)= 0$ for all $1\leq i<n$ we have $f(0)=f'(0)=0$ and follows from \eqref{conta-h-bordo} that $N(p)=\xi$. So,
  \[
  f''(0)= \langle X_{tt}, \xi \rangle (p)= \langle \bar\nabla_\nu \nu, N \rangle (p)=  h(\nu, \nu) >0
  \]
since $h(\nu, \nu) >0$,  by Proposition \ref{lema}. Thus, $f>0$ for  $0<t<\epsilon$, if $\epsilon>0$ is sufficiently small.
Hence, $M$ is locally contained in the halfspace $\Pi_+:=\{ p\in\mathbb{R}^{n+1} : \langle \xi, p\rangle >0\}$ along $\Sigma$.
\end{proof}

\end{document}
